% Part III -- Mathematics in 2026
% Cutoff: 2026-09-15, following the main handbook.

\finalchapter{16}{Prizes and Congresses in 2026}

\begin{fact}[상태 표현]
최근 수학에서는 status를 구분한다.
\begin{itemize}
\item \textbf{published}: scholarly publication에 게재.
\item \textbf{accepted}: journal acceptance가 공표되었지만 final issue는 아직일 수 있음.
\item \textbf{preprint}: 공개 manuscript, ordinary journal publication 전.
\item \textbf{announced}: authors/institution이 결과를 공개 발표; publication status는 별도.
\item \textbf{formally verified}: proof assistant가 formal proof term을 검사; ordinary peer review와 다름.
\item \textbf{scheduled}: 공식 발표된 미래 행사·임명.
\end{itemize}
numerical evidence, preprint, formal verification, published theorem을 서로 바꾸어 말하지 않는다.
\end{fact}

\begin{center}
\small
\renewcommand{\arraystretch}{1.03}
\begin{longtable}{@{}L{0.09\textwidth}L{0.28\textwidth}L{0.55\textwidth}@{}}
\toprule
연도 & 상·행사 & 핵심 연결\\
\midrule
\endfirsthead
\toprule 연도 & 상·행사 & 핵심 연결\\ \midrule \endhead
\bottomrule\endlastfoot
2025 & Abel -- Masaki Kashiwara & algebraic analysis, $D$-modules, crystal bases\\
2025 & Breakthrough -- Dennis Gaitsgory & geometric Langlands in characteristic zero\\
2025 & Shaw -- Kenji Fukaya & Fukaya category, symplectic topology, mirror symmetry\\
2025 & ICTP--IMU Ramanujan -- Claudio Muñoz & nonlinear dispersive PDE, soliton stability\\
2025 & Basic Science Lifetime Awards & Shigefumi Mori; George Lusztig; Robert Tarjan -- MMP; representation theory; algorithms/data structures\\
2026 & Abel -- Gerd Faltings & arithmetic geometry, Mordell conjecture, Diophantine geometry\\
2026 & ICM -- Philadelphia, Jul. 23--30 & Fields Medals and principal IMU prizes\\
2026 & Fields & Yu Deng; John Pardon; Jacob Tsimerman; Hong Wang\\
2026 & IMU Abacus -- Shayan Oveis Gharan & algorithms, combinatorial optimization\\
2026 & Gauss Prize -- Yurii Nesterov & convex optimization, accelerated first-order methods\\
2026 & Chern Medal -- Graeme Segal & topology, geometry, mathematical physics\\
2026 & Leelavati Award -- Hannah Fry & mathematics communication\\
2026 & Breakthrough -- Frank Merle & nonlinear evolution PDE, singularity, solitons\\
2026 & Shaw -- Emmanuel Candès; Camillo De Lellis & compressed sensing/statistics; GMT/fluid singularities\\
2026 & WLF Prize for Young Scientists -- Ziquan Zhuang & inaugural young-scientist prize; $K$-stability, birational geometry\\
2026 & Clay Research Awards & Furstenberg/Kakeya; Telescope counterexamples; Boltzmann from hard spheres\\
2026 & Rolf Schock -- Tobias Holck Colding & minimal surfaces, geometric flows\\
2026 & Maryam Mirzakhani Prize -- Roman Bezrukavnikov & geometric and modular representation theory\\
2026 & Michael and Sheila Held Prize & Irit Dinur; Subhash Khot; Guy Kindler; Dor Minzer; Muli Safra -- 2-to-2 Games, hardness, Unique Games\\
2026 & ICBS -- Beijing, Aug. 9--21 & inaugural ICBS mathematics medals\\
\end{longtable}
\end{center}

\begin{fact}[2026 Fields Medalists]
네 명의 이름과 대표 분야를 평행하게 기억한다.
\begin{itemize}
\item \textbf{Yu Deng}: PDE, kinetic theory.
\item \textbf{John Pardon}: symplectic geometry and topology.
\item \textbf{Jacob Tsimerman}: arithmetic geometry, o-minimality; 2026 season에는 AI-safety/MAISI 연결도 중요.
\item \textbf{Hong Wang}: harmonic analysis, geometric measure theory, Kakeya; third woman Fields Medalist.
\end{itemize}
\end{fact}

\begin{fact}[2026 ICBS 수학 메달]
2026 Beijing ICBS에서 inaugural mathematics medals:
\[
\text{Claire Voisin}\leftrightarrow\text{Emmy Noether Medal},
\]
\[
\text{Horng-Tzer Yau}\leftrightarrow\text{Shiing-Shen Chern Medal},
\]
\[
\boxed{\text{Shou-Wu Zhang}\leftrightarrow\text{Andrew Wiles Medal}}.
\]
각각 Hodge theory/algebraic geometry, random matrices/mathematical physics, arithmetic geometry가 강한 anchor.
\end{fact}

\finalchapter{17}{Mathematical Competitions in 2026}

\begin{center}
\small
\renewcommand{\arraystretch}{1.10}
\begin{tabularx}{0.98\textwidth}{@{}L{0.18\textwidth}L{0.26\textwidth}Y@{}}
\toprule
대회 & 대한민국 & 기억할 기록\\
\midrule
2025 IMO & 3위, 203점; 금 4, 은 2 & Sunshine Coast; gold cutoff $35$; perfect score 5명\\
2026 IMO & 6위, 167점; 금 3, 은 2, HM 1 & Shanghai; Hyeonjun Lee $42/42$; 세계 7명 만점\\
2025 APMO & Kyungjun Park $35/35$, 금메달 & 5문제, 각 7점\\
2026 APMO & 2위, 306점 & Seohyung Jo, Hyeonjun Lee, Hyunjun Jang 모두 $35/35$; 공식 award는 gold/silver/silver\\
\bottomrule
\end{tabularx}
\end{center}

\begin{fact}[2025·2026 IMO 대한민국 기록]
2025 대표 6명의 점수는
\[
37,36,36,35,31,28
\]
로 합계 $203$. 2026에는 Hyeonjun Lee가 $42/42$ perfect score를 기록했고 대한민국은 총점 $167$로 6위였다.
\end{fact}

\begin{fact}[2026 APMO]
대한민국은 43개국 중 2위, team total $306$. Seohyung Jo, Hyeonjun Lee, Hyunjun Jang이 모두 $35/35$였지만 공식 award label은 각각 Gold, Silver, Silver였다. 수치와 award label을 구분한다.
\end{fact}

\begin{fact}[2026 IMO Medal]
2026 IMO medal의 reverse에는 Zhao Shuang의 Pythagorean-theorem diagram이 사용되었고 지름은
\[
\boxed{9\,\mathrm{cm}}.
\]
현대 대회, 중국 수학사, 피타고라스 정리를 연결하는 강한 식별 단서이다.
\end{fact}

\finalchapter{18}{Major Mathematical Developments in 2026}

\begin{center}
\small
\renewcommand{\arraystretch}{1.08}
\begin{longtable}{@{}L{0.27\textwidth}L{0.22\textwidth}L{0.43\textwidth}@{}}
\toprule
주제 & 인물·상태 & 최종 암기 포인트\\
\midrule
\endfirsthead
\toprule 주제 & 인물·상태 & 최종 암기 포인트\\ \midrule \endhead
\bottomrule\endlastfoot
Navier--Stokes & OpenAI announcement; Clay evaluation pending at cutoff & forced 3D finite-time singularity manuscript + Lean formalization; Clay-recognized solution이라고 단정하지 않음\\
3D Kakeya & Hong Wang; Joshua Zahl; 2025 preprint, 2026 Clay Research Award & every Kakeya set in $\R^3$ has Hausdorff and Minkowski dimension $3$\\
Hilbert 10 over rings of integers & Alpöge--Bhargava--Ho--Shnidman; independently Koymans--Pagano & number fields의 integer rings에서도 general decision algorithm 없음\\
Boltzmann from hard spheres & Yu Deng; Zaher Hani; Xiao Ma & long-time rigorous derivation; Hilbert's sixth problem program\\
Ordinary abelian surfaces & Boxer--Calegari--Gee--Pilloni; 2025 preprint & modularity, Siegel modular forms, Galois representations\\
Mizohata--Takeuchi & Hannah Mira Cairo; counterexample preprint & Fourier restriction; logarithmic-loss phenomenon\\
Compact Bonnet pairs & Bobenko--Hoffmann--Sageman-Furnas; published 2025 & noncongruent immersed tori with same metric and mean curvature\\
Jacobian Conjecture & source-set status: counterexamples for all $n>2$; $n=2$ unresolved & dimension-dependent status가 핵심; source attribution 불일치 때문에 사람 이름은 Final에서 고정하지 않음\\
Riemann zeta critical line & Alpöge--Furman; independently Lamzouri & nontrivial zeros 중 $>2/3$가 simple and on $\Re(s)=1/2$; RH는 미해결\\
Sendov Conjecture & Lech Mazur; Tao exposition/Lean connection & 2026 AI-assisted proof; strengthened Phelps--Rodriguez form까지 연결\\
Erdős--Szemerédi sum-product over $\R$ & Bloom--Sawin--Schildkraut--Zhelezov & real-version counterexample; integer version과 구별\\
Petersen Coloring & Bryce Putman; Jorik Jooken & explicit 112-vertex counterexample\\
HRT Conjecture & Faulhuber et al. & time-frequency shifts의 linear independence conjecture가 false\\
Ravenel Telescope & Burklund--Hahn--Levy--Schlank & chromatic homotopy theory의 counterexamples; 2026 Clay Research Award\\
\end{longtable}
\end{center}

\begin{fact}[Navier--Stokes Status Trap]
안전한 표현은
\[
\boxed{\text{해결 발표·원고/형식화 공개; Clay 평가 진행 중}}.
\]
formal verification, ordinary publication, Clay recognition, credit determination을 서로 같은 상태로 취급하지 않는다.
\end{fact}

\begin{fact}[Jacobian Conjecture Status]
과학퀴즈 source set에서 외울 핵심은
\[
n>2:\ \text{counterexamples},\qquad n=2:\ \text{unresolved}.
\]
서로 다른 source에서 attribution이 일치하지 않았으므로 Final에서는 사람 이름보다 status를 우선한다.
\end{fact}

\begin{fact}[Zeta Numerical Anchor]
\[
\boxed{>\frac23}
\]
의 nontrivial zeros가 simple하면서 critical line 위에 있다는 unconditional 2026 lower-bound result. ``RH 해결''이 아니다.
\end{fact}

\finalchapter{19}{Artificial Intelligence and Mathematics in 2026}

\begin{fact}[2025 IMO AI: Gold-Medal-Level]
Google DeepMind의 Gemini Deep Think는 2025 IMO 6문제 중 5문제를 풀어 $35/42$, human gold cutoff와 같은 점수를 기록했다고 보고되었다.
OpenAI도 별도의 general-purpose reasoning model이 $35/42$를 기록했다고 보고했다.
둘 다 공식 참가자가 아니므로 ``gold-medal-level''이지 실제 IMO medalist는 아니다.
\end{fact}

\begin{fact}[AlphaProof]
2025 Nature paper. reinforcement learning + Lean proof assistant.
2024 IMO benchmark에서 geometry system과 결합해 silver-medal-level performance가 알려졌다.
핵심은 \textbf{machine-checkable formal proofs}.
\end{fact}

\begin{fact}[AlphaEvolve]
Google DeepMind가 2025 공개. evolutionary coding agent로 algorithm/construction space를 탐색하고 executable candidates를 자동 평가.
단순 prose generation과 구별한다.
\end{fact}

\begin{fact}[GPT-5.5 and Off-Diagonal Ramsey Numbers]
2026년 4월 OpenAI는 internal GPT-5.5 system이 off-diagonal Ramsey numbers의 오래된 asymptotic fact에 대한 새 proof를 찾았고
이후 Lean verification이 이루어졌다고 보고했다.
\[
\text{GPT-5.5}\leftrightarrow\text{off-diagonal Ramsey}\leftrightarrow\text{Lean}.
\]
\end{fact}

\begin{fact}[GPT-5.6 Sol and Cycle Double Cover]
Cycle Double Cover Conjecture는 every finite bridgeless loopless multigraph의 각 edge가 정확히 두 번 포함되는 cycle multiset의 존재를 묻는다.
2026년 7월 OpenAI는 GPT-5.6 Sol Ultra가 생성한 affirmative proof 자료를 공개했고,
Sang-il Oum과 Jim Geelen이 독립 expository write-ups를 게시했다.
ordinary journal publication과는 구분한다.
\end{fact}

\begin{fact}[First Proof]
2026년 2월 OpenAI가 research-level First Proof challenge 10문제 전체에 대한 proof attempts를 공개.
일부는 높은 correct-probability 평가를 받았지만 적어도 하나의 초기 favored attempt는 이후 incorrect로 인정되었다.
핵심 lesson: convincing research proof attempt $\neq$ independently checked theorem.
\end{fact}

\begin{fact}[AI Disproof of the Erdős Unit-Distance Conjecture]
2026년 5월 공개된 AI-generated argument는 infinitely many $n$에 대해 어떤 fixed $\delta>0$로
\[
u(n)\ge n^{1+\delta}
\]
를 주어 Erdős의 near-linear conjecture를 반박한다.
exact asymptotic order를 결정한 결과는 아니다.
\end{fact}

\begin{fact}[Palomar]
2026년 8월 Jeremy Avigad, Terence Tao, Ravi Vakil, Akshay Venkatesh 등이 포함된 group이
Lean-verified mathematics registry \textbf{Palomar}를 발표.
fixed source versions, exact formal statements, dependencies, review comments를 기록한다.
kernel checking은 novelty나 ordinary peer review를 자동으로 보증하지 않는다.
\end{fact}

\begin{fact}[Fermat's Last Theorem in Lean]
2026년 9월 4일 Anthropic은 Claude와 Prove2Me로 만든 FLT formalization을 공개.
Darmon--Diamond--Taylor exposition을 따라 Wiles의 argument를 Lean으로 computer-check.
\textbf{새로운 FLT proof가 아니라 formal-verification milestone}.
\end{fact}

\begin{fact}[Prime Gaps at Most 186]
2026년 9월 PrimeGaps186 project는
\[
\liminf_{n\to\infty}(p_{n+1}-p_n)\le186
\]
의 Lean derivation을 다루지만 explicit input axioms와 numerical certificate에 의존하는 conditional formalization.
Twin Prime Conjecture가 해결된 것은 아니다.
\end{fact}

\begin{fact}[Leiden Declaration]
2026 Leiden Declaration은 AI-assisted mathematics에서
transparent disclosure, accurate attribution, methods/evidence에 대한 openness, human responsibility를 강조.
IMU endorsement가 season-specific anchor.
\end{fact}

\begin{strategy}
AI 관련 문제에서는
\[
\text{score},\quad \text{discovery},\quad \text{formal verification},\quad
\text{peer review},\quad \text{publication}
\]
을 서로 다른 주장으로 본다.
\end{strategy}


\begin{fact}[Ten AI-Generated Research Advances]
2026년 8월 OpenAI는 internal Astra 계열 model이 생성하고 Lean certificate까지 붙인 10개 research-result manuscripts를 공개했다. 대표 주제는 sphere packing, binary/spherical codes, non-sofic groups, Connes rigidity, arithmetic-circuit lower bounds, quantum parallel repetition, closest vector hardness, Ehrhart volume, multicolor Ramsey, extremal graph theory. cutoff에서는 \textbf{recently released manuscripts + formal certificates}로 기억한다.
\end{fact}

\begin{fact}[Claude and the Zeta Critical-Line Record]
Alpöge--Furman의 2026 paper는 Claude가 mathematical argument를 발견·작성했다고 명시하고, listed authors가 proof를 검증하고 책임을 진다고 설명한다. Lamzouri가 독립적으로 같은 bound의 더 짧은 proof를 게시. 수학적 내용은 Chapter 18의 $>2/3$ critical-line record와 연결한다.
\end{fact}

\finalchapter{20}{Mathematicians and Institutions in 2026}

\begin{center}
\small
\begin{longtable}{@{}L{0.26\textwidth}L{0.18\textwidth}L{0.48\textwidth}@{}}
\toprule
인물 & 2025--2026 event & 강한 연결\\
\midrule
\endfirsthead
\toprule 인물 & event & 강한 연결\\ \midrule \endhead
\bottomrule\endlastfoot
Peter D. Lax & died May 16, 2025, aged 99 & PDE, conservation laws, numerical analysis; 1987 Wolf, 2005 Abel\\
Heisuke Hironaka & died Mar. 18, 2026, aged 94 & resolution of singularities; 1970 Fields\\
Michael O. Rabin & died Apr. 14, 2026 & finite automata, probabilistic algorithms; 1976 Turing Award\\
Cleve Moler & died May 20, 2026, aged 86 & creator of MATLAB, numerical linear algebra; cofounder of MathWorks\\
James Maynard & Oxford Regius appointment & successor to Andrew Wiles, scheduled Oct. 1, 2026\\
Jacob Tsimerman & AI-safety connection & University of Toronto professor on leave; MAISI Scientific Director\\
\bottomrule
\end{longtable}
\end{center}

\begin{fact}[Cleve Moler]
가장 짧은 퀴즈 연결:
\[
\boxed{\text{Cleve Moler}\leftrightarrow\text{MATLAB}\leftrightarrow\text{numerical computing}}.
\]
\end{fact}

\begin{fact}[Wiles에서 Maynard로]
Andrew Wiles는 Oxford의 modern Regius Professorship of Mathematics의 inaugural holder.
James Maynard가 successor로 임명되어 2026년 10월 1일 시작 예정.
\end{fact}

\finalchapter{21}{Mathematics in the Republic of Korea in 2026}

\begin{center}
\small
\begin{longtable}{@{}L{0.22\textwidth}L{0.28\textwidth}L{0.42\textwidth}@{}}
\toprule
인물·행사 & 한국 연결 & 강한 퀴즈 단서\\
\midrule
\endfirsthead
\toprule 인물·행사 & 한국 연결 & 강한 퀴즈 단서\\ \midrule \endhead
\bottomrule\endlastfoot
Sug Woo Shin & KIAS distinguished professor; UC Berkeley & 2025 Samsung Ho-Am; number theory, automorphic forms, Langlands\\
Jae Kyoung Kim & IBS Biomedical Mathematics & 2025 SIAM Annual Meeting first Korean plenary speaker; mathematical biology\\
Nam-Gyu Kang & KIAS & 2025 KMS DI Prize; probability, complex analysis, CFT\\
Inhyeok Choi & KIAS & 2025 Sangsan Prize; geometric group theory, random walks\\
Jae Choon Cha & \POSTECH; Ph.D. from \KAIST & 2025 Korea Science Award; geometric topology, Disk Embedding Theory, Milnor problem\\
Cheolwoo Park & \KAIST & 2025 Choi Seok-jeong Award; statistical machine learning\\
Sung-Jin Oh & \KAIST alumnus; KIAS/UC Berkeley & 2026 Samsung Ho-Am; nonlinear hyperbolic PDE, general relativity\\
Bo-Hae Im & \KAIST & 2026 KMS Paper Award; positive-characteristic multiple zeta values\\
JungHwan Park & \KAIST & 2026 KMS Paper Award; knot concordance, $(2,1)$-cable of figure-eight knot\\
POSTECH Arithmetic Geometry Lab & \POSTECH & 2026 Basic Research Laboratory; Langlands, Diophantine/arithmetic geometry\\
Kyeongsu Choi & KIAS & 2026 POSCO TJ Park Science Prize; curvature flows, geometric PDE\\
IMU Executive Committee & Seoul, KIAS, Mar. 13--14 & IMU and Korean mathematical institutions\\
ASIACOMB & Daejeon, Aug. 24--28 & inaugural Asian combinatorics forum; Huy Tuan Pham inaugural prize\\
Jineon Baek & \POSTECH alumnus; June E Huh Fellow at KIAS & moving sofa, Gerver area $\approx2.21953166$, preprint status\\
\bottomrule
\end{longtable}
\end{center}

\begin{fact}[차재춘 (Jae Choon Cha)]
\POSTECH{} professor, Ph.D. from \KAIST.
\[
\text{Jae Choon Cha}\leftrightarrow\text{geometric topology}\leftrightarrow\text{4D Disk Embedding Theory / transfinite invariants}.
\]
2025 Korea Science Award와 두 학교를 동시에 잇는 점이 Science War에서 특히 강한 단서.
\end{fact}

\begin{fact}[KAIST 2026 KMS 논문상]
\textbf{Bo-Hae Im}: positive-characteristic multiple/alternating multiple zeta values, Zagier--Hoffman-type conjectures.
\textbf{JungHwan Park}: knot concordance; $(2,1)$-cable of the figure-eight knot가 smoothly slice가 아님.
\end{fact}
